\begin{ouverture}[grandeurs]{9}{Proportionnalité et échelles}{Partie 2 · Modéliser et mesurer}
{Grandeurs · Linéarité · Retour à l’unité · Échelles}
\objectif{Reconnaître une situation de proportionnalité entre deux grandeurs.}
\objectif{Choisir une procédure adaptée : fois plus, sommes, retour à l’unité.}
\objectif{Représenter une situation dans un tableau et utiliser une échelle.}
\prerequis{Multiplication et division ; « fois plus, fois moins » du cours moyen.}
\end{ouverture}
\pagestyle{otMain}\essentiel\fontsize{11}{15}\selectfont
\phantomsection\label{cours-9-1}
\essSection{1}{Reconnaître une situation proportionnelle}
\begin{definition}Deux grandeurs sont proportionnelles lorsque, si l’on multiplie (ou divise) une mesure de l’une par un nombre, la mesure correspondante de l’autre est multipliée (ou divisée) par le même nombre. Exemple : la masse de pommes achetées et le prix payé, à $2{,}50$ € le kilogramme.\end{definition}
\begin{propriete}Dans une situation de proportionnalité, une grandeur « par unité » reste constante : le prix au kilogramme, le nombre de battements du cœur par minute, la quantité de farine par personne.\end{propriete}
\begin{attention}Toutes les situations ne sont pas proportionnelles : la taille d’un enfant ne double pas quand son âge double ; une course de taxi avec une prise en charge ne coûte pas deux fois plus cher pour deux fois plus de kilomètres.\end{attention}
\essSection{2}{Utiliser les relations entre les mesures}
\begin{savoirfaire}[Fois plus, fois moins]4 croissants coûtent $3{,}60$ €. 12 croissants, c’est 3 fois plus : ils coûtent 3 fois plus, $10{,}80$ €. 2 croissants, c’est 2 fois moins : $1{,}80$ €.\end{savoirfaire}
\begin{savoirfaire}[Additionner deux situations]2 kg de farine coûtent $3{,}20$ € et 3 kg coûtent $4{,}80$ €. Donc 5 kg, soit 2 kg et 3 kg, coûtent $3{,}20+4{,}80=8$ €.\end{savoirfaire}
\clearpage
\phantomsection\label{cours-9-2}
\essSection{3}{Revenir à l’unité et représenter}
\begin{savoirfaire}[Retour à l’unité]5 cahiers coûtent $7{,}50$ €. Un cahier coûte $7{,}50:5=1{,}50$ € ; 8 cahiers coûtent donc $8\times1{,}50=12$ €.\end{savoirfaire}
\label{fig-cours-9-2-cours-tableau-fleches}\figureGeo{cours-tableau-fleches}

\begin{exemple}Dans un tableau, on écrit le nom de chaque grandeur et son unité. On peut aussi noter : prix(2 kg) $=5$ € ; prix(6 kg) $=3\times$ prix(2 kg) $=15$ €.\end{exemple}
\essSection{4}{Lire et utiliser une échelle}
\begin{definition}Sur un plan à l’échelle, les longueurs du plan et les longueurs réelles sont proportionnelles. « 1 cm pour 50 m » signifie que 1 cm sur le plan représente 50 m en réalité.\end{definition}
\label{fig-cours-9-2-cours-echelle}\figureGeo{cours-echelle}

\begin{savoirfaire}[Du plan à la réalité]6,5 cm sur le plan représentent $6{,}5\times50=325$ m. Une distance réelle de 400 m mesure $400:50=8$ cm sur le plan.\end{savoirfaire}

\clearpage\rubrique{Découverte}
\begin{blocLivre}{ex-9-1}\begin{exercice}[1]{Proportionnel ou pas ?}
Dire si les deux grandeurs sont proportionnelles, en justifiant.\begin{questions}
\item Le prix payé et le nombre de baguettes achetées, à $1{,}20$ € la baguette.
\item La taille d’un enfant et son âge.
\item Le périmètre d’un carré et la longueur de son côté.
\item Le prix d’une course de taxi (4 € de prise en charge et 2 € par kilomètre) et la distance parcourue.
\end{questions}
\end{exercice}\end{blocLivre}
\begin{blocLivre}{ex-9-2}\begin{exercice}[1]{Trois fois plus}
4 croissants coûtent $3{,}60$ €. Combien coûtent 12 croissants ? Et 2 croissants ?
\end{exercice}\end{blocLivre}
\begin{blocLivre}{ex-9-3}\begin{exercice}[1]{Retour à l’unité}
5 cahiers identiques coûtent $7{,}50$ €. Combien coûtent 8 cahiers ?
\end{exercice}\end{blocLivre}
\begin{blocLivre}{ex-9-4}\begin{exercice}[1]{Additionner des situations}
2 kg de farine coûtent $3{,}20$ € et 3 kg coûtent $4{,}80$ €. Combien coûtent 5 kg ? Et 1 kg ?
\end{exercice}\end{blocLivre}
\begin{blocLivre}{ex-9-5}\begin{exercice}[1]{Lire une échelle}
Sur le plan, 1 cm représente 200 m. La distance entre l’école et le gymnase mesure 4 cm sur le plan. Quelle est la distance réelle ?
\label{fig-ex-9-5}\figureGeo{pro-5}
\end{exercice}\end{blocLivre}

\rubrique{Entraînement}
\begin{blocLivre}{ex-9-6}\begin{exercice}[2]{Compléter un tableau}
Le prix d’un fromage est proportionnel à sa masse. Recopier et compléter le tableau.\begin{center}\renewcommand{\arraystretch}{1.35}\begin{tabular}{c|cccc}\toprule
Masse (kg) & 2 & 3 & 5 & 10\\\midrule
Prix (€) & 7 & … & … & …\\
\bottomrule\end{tabular}\end{center}
\end{exercice}\end{blocLivre}
\begin{blocLivre}{ex-9-7}\begin{exercice}[2]{Les battements du cœur}
Au repos, le cœur de Sam bat régulièrement 72 fois par minute. Combien de battements en 5 minutes ? En 30 secondes ? En une heure ?
\end{exercice}\end{blocLivre}
\begin{blocLivre}{ex-9-8}\begin{exercice}[2]{Proportionnel ?}
Un parc d’attractions affiche ses tarifs. Le prix est-il proportionnel au nombre de places ?\begin{center}\renewcommand{\arraystretch}{1.35}\begin{tabular}{c|cccc}\toprule
Nombre de places & 1 & 2 & 3 & 4\\\midrule
Prix (€) & 8 & 16 & 24 & 30\\
\bottomrule\end{tabular}\end{center}
\end{exercice}\end{blocLivre}
\begin{blocLivre}{ex-9-9}\begin{exercice}[2]{Le plan du jardin}
Sur ce plan, 1 cm représente 2 m. Le jardin est représenté par un rectangle de 6 cm sur $4{,}5$ cm. Quelles sont ses dimensions réelles ?
\label{fig-ex-9-9}\figureGeo{pro-9}
\end{exercice}\end{blocLivre}
\begin{blocLivre}{ex-9-10}\begin{exercice}[2]{La recette}
Pour 6 personnes, une recette demande 300 g de farine, 3 œufs et $0{,}6$ L de lait. Adapter la recette pour 4 personnes.
\end{exercice}\end{blocLivre}

\rubrique{Problèmes}
\begin{blocLivre}{ex-9-11}\begin{exercice}[3]{La carte de randonnée}
Sur la carte, 1 cm représente 250 m. Le chemin suivi est formé de trois portions de 2 cm, $3{,}4$ cm et $1{,}6$ cm. Quelle distance réelle parcourt-on ?
\label{fig-ex-9-11}\figureGeo{pro-11}
\end{exercice}\end{blocLivre}
\begin{blocLivre}{ex-9-12}\begin{exercice}[3]{Le robinet}
Un robinet débite régulièrement 12 L d’eau en $1{,}5$ min. Quelle quantité en 6 min ? En 1 min ? Combien de temps faut-il pour remplir une cuve de 100 L ?
\end{exercice}\end{blocLivre}
\begin{blocLivre}{ex-9-13}\begin{exercice}[3]{Le meilleur prix}
Un paquet de café de 500 g coûte $2{,}40$ € ; un paquet de 750 g coûte $3{,}30$ €. Calculer le prix au kilogramme de chaque paquet. Lequel est le plus avantageux ?
\end{exercice}\end{blocLivre}
\begin{blocLivre}{ex-9-14}\begin{exercice}[3]{La maquette}
Sur une maquette d’avion, 1 cm représente 50 cm de l’avion réel. L’envergure de la maquette est de 24 cm : quelle est celle de l’avion ? L’avion mesure 15 m de long : quelle est la longueur de la maquette ?
\end{exercice}\end{blocLivre}
\begin{blocLivre}{ex-9-15}\begin{exercice}[3]{Pas proportionnel !}
Un taxi facture 4 € de prise en charge, puis 2 € par kilomètre. Calculer le prix d’une course de 5 km, puis de 10 km. Le prix est-il proportionnel à la distance ? Quelle partie du prix l’est ?
\end{exercice}\end{blocLivre}

\clearpage
